% GENERATED FILE. Edit canonical lessons, then run npm run generate:books.
% canonical-source-hash: fnv1a64:954452b28e51f330
% canonical-lessons: BN-C365-motbhed, BN-C365-somjhota, BN-C365-nishpotti, BN-C365-manano, BN-C365-sorbosommoto

\chapter{Coming to Terms}
\label{ch:bn-coming-to-terms}
\hlchaptermodality{365}

\begin{chapteropening}
\textbf{By the end of this chapter:} \emph{I can say how a disagreement was settled.}

Complete the last lesson of chapter 365: I can say how a disagreement was settled.
\end{chapteropening}

\section[\texorpdfstring{\bn{মতভেদ} (môtbhed)}{মতভেদ (môtbhed)}]{\bn{মতভেদ} (môtbhed) — a difference of opinion}

\label{lesson:BN-C365-motbhed}



\begin{quote}
Before the new one: say the Bengali for to be devastated, then the Bengali for extreme.
\end{quote}

\subsection*{You'll want to know: \bn{মতভেদ}}
\textbf{\bn{মতভেদ}} — \emph{môtbhed} — \textquotedblleft{}a difference of opinion\textquotedblright{}.

\begin{quote}
\bn{এই} \bn{বিষয়ে} \bn{আমাদের} \bn{মতভেদ} \bn{আছে।}
\end{quote}

\begin{quote}
\emph{ei bishôye āmāder môtbhed āchhe.}
\end{quote}

\begin{quote}
We disagree on this subject.
\end{quote}

\begin{cousinweb}
\textbf{\bn{মতভেদ}} (\emph{môtbhed}) is Sanskrit \emph{mata}, an opinion, with \emph{bheda}, a split or difference (Monier-Williams): where opinions part.
\end{cousinweb}

\begin{grammarlens}[title={What you've built}]
One more word for giving reasons and opinions.
\end{grammarlens}

\subsection*{Your turn}
\begin{itemize}
\raggedright
  \item \emph{Say it:} \emph{môtbhed}
  \item \emph{Say it:} \emph{môtbhed}, once more
  \item \emph{Say it:} \emph{ei bishôye āmāder môtbhed āchhe.} — We disagree on this subject.
  \item \emph{Recall:} say the Bengali for to be devastated, then the Bengali for extreme, then say \emph{môtbhed} again
\end{itemize}

\begin{tcolorbox}[breakable,colback=teal!4,colframe=teal!35!black,title={Before you move on}]
What does \bn{মতভেদ} mean? (\textquotedblleft{}A difference of opinion\textquotedblright{}.) Say it once more.
\end{tcolorbox}
\section[\texorpdfstring{\bn{সমঝোতা} (sômjhotā)}{সমঝোতা (sômjhotā)}]{\bn{সমঝোতা} (sômjhotā) — an agreement reached; an understanding}

\label{lesson:BN-C365-somjhota}



\begin{quote}
Before the new one: say the Bengali for extreme, then the Bengali for a difference of opinion.
\end{quote}

\subsection*{You'll want to know: \bn{সমঝোতা}}
\textbf{\bn{সমঝোতা}} — \emph{sômjhotā} — \textquotedblleft{}an agreement reached; an understanding\textquotedblright{}.

\begin{quote}
\bn{শেষে} \bn{দুই} \bn{দল} \bn{সমঝোতায়} \bn{এল।}
\end{quote}

\begin{quote}
\emph{sheshe dui dôl sômjhotāy elô.}
\end{quote}

\begin{quote}
In the end the two sides reached an agreement.
\end{quote}

\begin{cousinweb}
\textbf{\bn{সমঝোতা}} (\emph{sômjhotā}) is usually traced to Hindustani \emph{samjhautā}, an understanding, built on \emph{samajh}, understanding: what two sides reach when each understands the other. \bn{সমঝোতায়} \bn{আসা} is to come to terms.
\end{cousinweb}

\begin{grammarlens}[title={What you've built}]
One more word for giving reasons and opinions.
\end{grammarlens}

\subsection*{Your turn}
\begin{itemize}
\raggedright
  \item \emph{Say it:} \emph{sômjhotā}
  \item \emph{Say it:} \emph{sômjhotā}, once more
  \item \emph{Say it:} \emph{sheshe dui dôl sômjhotāy elô.} — In the end the two sides reached an agreement.
  \item \emph{Recall:} say the Bengali for extreme, then the Bengali for a difference of opinion, then say \emph{sômjhotā} again
\end{itemize}

\begin{tcolorbox}[breakable,colback=teal!4,colframe=teal!35!black,title={Before you move on}]
What does \bn{সমঝোতা} mean? (\textquotedblleft{}An agreement reached; an understanding\textquotedblright{}.) Say it once more.
\end{tcolorbox}
\section[\texorpdfstring{\bn{নিষ্পত্তি} (nishpôtti)}{নিষ্পত্তি (nishpôtti)}]{\bn{নিষ্পত্তি} (nishpôtti) — a settlement; a resolution (of a case or a dispute)}

\label{lesson:BN-C365-nishpotti}



\begin{quote}
Before the new one: say the Bengali for a difference of opinion, then the Bengali for an agreement reached.
\end{quote}

\subsection*{You'll want to know: \bn{নিষ্পত্তি}}
\textbf{\bn{নিষ্পত্তি}} — \emph{nishpôtti} — \textquotedblleft{}a settlement; a resolution (of a case or a dispute)\textquotedblright{}.

\begin{quote}
\bn{মামলাটার} \bn{নিষ্পত্তি} \bn{হয়ে} \bn{গেল।}
\end{quote}

\begin{quote}
\emph{māmlāṭār nishpôtti hôye gelô.}
\end{quote}

\begin{quote}
The case was settled.
\end{quote}

\begin{cousinweb}
\textbf{\bn{নিষ্পত্তি}} (\emph{nishpôtti}) is Sanskrit \emph{niṣpatti}, completion, a thing brought to its end (Monier-Williams). Bengali uses it for the final settling of a case or a quarrel.
\end{cousinweb}

\begin{grammarlens}[title={What you've built}]
One more word for giving reasons and opinions.
\end{grammarlens}

\subsection*{Your turn}
\begin{itemize}
\raggedright
  \item \emph{Say it:} \emph{nishpôtti}
  \item \emph{Say it:} \emph{nishpôtti}, once more
  \item \emph{Say it:} \emph{māmlāṭār nishpôtti hôye gelô.} — The case was settled.
  \item \emph{Recall:} say the Bengali for a difference of opinion, then the Bengali for an agreement reached, then say \emph{nishpôtti} again
\end{itemize}

\begin{tcolorbox}[breakable,colback=teal!4,colframe=teal!35!black,title={Before you move on}]
What does \bn{নিষ্পত্তি} mean? (\textquotedblleft{}A settlement; a resolution (of a case or a dispute)\textquotedblright{}.) Say it once more.
\end{tcolorbox}
\section[\texorpdfstring{\bn{মানানো} (mānāno)}{মানানো (mānāno)}]{\bn{মানানো} (mānāno) — to talk round; to win over; to suit}

\label{lesson:BN-C365-manano}



\begin{quote}
Before the new one: say the Bengali for an agreement reached, then the Bengali for a settlement.
\end{quote}

\subsection*{You'll want to know: \bn{মানানো}}
\textbf{\bn{মানানো}} — \emph{mānāno} — \textquotedblleft{}to talk round; to win over; to suit\textquotedblright{}.

\begin{quote}
\bn{শেষে} \bn{আমি} \bn{বাবাকে} \bn{মানালাম।}
\end{quote}

\begin{quote}
\emph{sheshe āmi bābāke mānālām.}
\end{quote}

\begin{quote}
In the end I talked my father round.
\end{quote}

Past (he/she …): \bn{তিনি} \bn{সবাইকে} \bn{মানালেন}, \emph{tini sôbāike mānālen}.

\begin{cousinweb}
\textbf{\bn{মানানো}} (\emph{mānāno}) is the causative of \bn{মানা} 'to accept' (CDIAL 9857 lists Bengali mānā): to make someone accept, to bring them round. A colour or a dress \bn{মানায়} when it suits the one wearing it.
\end{cousinweb}

\begin{grammarlens}[title={What you've built}]
One more word for giving reasons and opinions.
\end{grammarlens}

\subsection*{Your turn}
\begin{itemize}
\raggedright
  \item \emph{Say it:} \emph{mānāno}
  \item \emph{Say it:} \emph{mānāno}, once more
  \item \emph{Say it:} \emph{sheshe āmi bābāke mānālām.} — In the end I talked my father round.
  \item \emph{Recall:} say the Bengali for an agreement reached, then the Bengali for a settlement, then say \emph{mānāno} again
\end{itemize}

\begin{tcolorbox}[breakable,colback=teal!4,colframe=teal!35!black,title={Before you move on}]
What does \bn{মানানো} mean? (\textquotedblleft{}To talk round; to win over; to suit\textquotedblright{}.) Say it once more.
\end{tcolorbox}
\section[\texorpdfstring{\bn{সর্বসম্মত} (sôrbôsômmôtô)}{সর্বসম্মত (sôrbôsômmôtô)}]{\bn{সর্বসম্মত} (sôrbôsômmôtô) — unanimous}

\label{lesson:BN-C365-sorbosommoto}



\begin{quote}
Before the new one: say the Bengali for a settlement, then the Bengali for to talk round.
\end{quote}

\subsection*{You'll want to know: \bn{সর্বসম্মত}}
\textbf{\bn{সর্বসম্মত}} — \emph{sôrbôsômmôtô} — \textquotedblleft{}unanimous\textquotedblright{}.

\begin{quote}
\bn{সিদ্ধান্তটা} \bn{সর্বসম্মত} \bn{ছিল।}
\end{quote}

\begin{quote}
\emph{siddhāntôṭā sôrbôsômmôtô chhilô.}
\end{quote}

\begin{quote}
The decision was unanimous.
\end{quote}

\begin{cousinweb}
\textbf{\bn{সর্বসম্মত}} (\emph{sôrbôsômmôtô}) is Sanskrit \emph{sarva}, all, with \emph{sammata}, agreed or approved (Monier-Williams): agreed by everyone.
\end{cousinweb}

\begin{grammarlens}[title={What you've built}]
One more word for giving reasons and opinions.
\end{grammarlens}

\subsection*{Your turn}
\begin{itemize}
\raggedright
  \item \emph{Say it:} \emph{sôrbôsômmôtô}
  \item \emph{Say it:} \emph{sôrbôsômmôtô}, once more
  \item \emph{Say it:} \emph{siddhāntôṭā sôrbôsômmôtô chhilô.} — The decision was unanimous.
  \item \emph{Recall:} say the Bengali for a settlement, then the Bengali for to talk round, then say \emph{sôrbôsômmôtô} again
\end{itemize}

\begin{tcolorbox}[breakable,colback=teal!4,colframe=teal!35!black,title={Before you move on}]
What does \bn{সর্বসম্মত} mean? (\textquotedblleft{}Unanimous\textquotedblright{}.) Say it once more.
\end{tcolorbox}
